\appendixsection{Полная таблица угроз}
\label{app:threats}

\begin{table}[htbp]
\caption{Реестр сценариев для проектируемой системы}
\centering\small
\begin{tabularx}{\linewidth}{|p{.08\linewidth}|X|X|}
\hline
Код & Источник, условие и объект & Последствие \\ \hline
У1 & Совместный доступ к состояниям двух узлов & Восстановление признаков или параметров дерева \\ \hline
У2 & Наблюдение ветвления по адресу следующего узла & Раскрытие исхода сравнения \\ \hline
У3 & Подмена узла или смешение версий при передаче долей & Раскрытие доли либо ошибочный прогноз \\ \hline
У4 & Получение результата посторонней стороной или накопление разрешённых ответов & Раскрытие метки или сведений о модели \\ \hline
У5 & Потеря вычислительного узла во время запроса & Нарушение доступности \\ \hline
У6 & Недостоверная модель или искажённый вход & Нарушение целостности решения \\ \hline
\end{tabularx}
\end{table}

\appendixsection{Конфигурация окружения пробного запуска}
\label{app:env}

Дата и идентификатор сборки, ОС, процессор, память, версии Python, TensorFlow и TF~Encrypted, режим транспорта, тайм-ауты и расположение журнала заполняются по фактическому стенду.
